% TOP_PROBLEM: {"id":30007624,"problem_number":"LOCAL-30007624","title":"Dynamic tariff incidence","category_id":21,"category":{"id":21,"name":"economics","display_name":"Economics","description":"Open economic theory statements, published research agendas, and proposed scoped specifications; question provenance and historical priority are qualified per record.","slug":"economics","order_index":21,"created_at":"2026-10-08T00:00:00Z"},"status":"research_question","external_url":null,"legacy_ids":[],"published":false,"statement":"In the explicitly specified setting: How does the burden of large tariffs shift over time among exporters, importers, retailers, workers and consumers? The mathematical statement above fixes the target, assumptions and scope of this catalogue formulation.","economics":{"original_id":113,"field":"Trade, globalization and geoeconomics","kind":"identification","formalization_basis":"adapted_research_question","source_status":"research_agenda","priority_status":"earliest_found","priority_note":"Earliest explicit uncertainty discussion located in this audit, not the historical first statement of tariff-incidence research. The March 2026 conference manuscript predates its April NBER issue.","earliest_verified_reference":{"authors":["Pablo D. Fajgelbaum","Amit Khandelwal"],"title":"Tariffs in 2025: Short-Run Impacts on the U.S. Economy","year":2026,"url":"https://www.brookings.edu/wp-content/uploads/2026/03/1_Fajgelbaum-Khandelwal_unembargoed.pdf","doi":"10.3386/w35064","locator":"Abstract, PDF page 2; Section 5.2.2, final paragraph, PDF page 30; Section 6.4, PDF page 44","evidence_summary":"Finds terms-of-trade evidence inconclusive; longer adjustment and manufacturing outcomes remain uncertain."},"current_reference":{"authors":["Pablo D. Fajgelbaum","Amit Khandelwal"],"title":"Tariffs in 2025: Short-Run Impacts on the U.S. Economy","year":2026,"url":"https://www.brookings.edu/wp-content/uploads/2026/03/1_Fajgelbaum-Khandelwal_unembargoed.pdf","doi":"10.3386/w35064","locator":"Abstract, PDF page 2; Section 5.2.2, final paragraph, PDF page 30; Section 6.4, PDF page 44","evidence_summary":"Finds terms-of-trade evidence inconclusive; longer adjustment and manufacturing outcomes remain uncertain."},"formalization_note":"The source explicitly leaves aggregate price adjustment uncertain and suggests it may require more time. The whole-chain dynamic pass-through estimands, inventories and welfare accounting are proposed extensions; many short-run subcases are already estimated."},"statusQualification":"Published question or research agenda verified at the cited locator. The mathematical specification is adapted for this catalogue and includes additional assumptions; source publication does not establish first-ever priority or certify this exact model remains unsolved. The source explicitly leaves aggregate price adjustment uncertain and suggests it may require more time. The whole-chain dynamic pass-through estimands, inventories and welfare accounting are proposed extensions; many short-run subcases are already estimated.","statusReviewedAt":"2026-10-08"}
% FIRST_POSTED: 2026-10-08
% ENTRY_KIND: statement_only
% !TeX program = lualatex
\documentclass[11pt]{article}
\usepackage{fontspec}
\usepackage[a4paper,margin=25mm]{geometry}
\usepackage{amsmath,amssymb,mathtools}
\usepackage{xurl}
\usepackage[hidelinks]{hyperref}
\newcommand{\sref}[2]{\hyperlink{source:#1:#2}{[S#2]}}
\setlength{\emergencystretch}{3em}
\begin{document}
% problemId:30007624
\section*{Dynamic tariff incidence}\label{30007624}
\subsection{Definitions and mathematical statement}
Fix a product-level tariff intervention increasing the statutory ad valorem rate from \(\tau_0\) to \(\tau_1\), with \(0\leq\tau_0<\tau_1\), at date \(t\). Let \(P_{k,t}\) denote the price at supply-chain stage \(k\) and date \(t\). Track foreign export prices, tariff-inclusive import prices, domestic producer prices, retail prices, quantities, wages and profits. For each stage \(k\) and horizon \(h\), define pass-through
\[\rho_k(h)=\frac{E[\log P_{k,t+h}(do(\tau_1))-\log P_{k,t+h}(do(\tau_0))]}{\log(1+\tau_1)-\log(1+\tau_0)}.\]
The task is to identify these dynamic responses and allocate the welfare burden across foreign producers, importers, retailers, workers and consumers, counting tariff revenue consistently. Specify product composition, exchange-rate treatment and the counterfactual policy path, including retaliation. Endogenous sourcing changes mean price indexes must distinguish within-product repricing from entry, exit and quality substitution. Anticipatory stockpiling and inventory drawdown must be modeled or measured before interpreting delayed retail effects. Establish conditions for a valid comparison group when tariff assignments target politically or economically distinct products. Report identified sets where export margins or profits are unobserved. High import-price pass-through and low contemporaneous consumer-price pass-through are compatible and do not resolve how burdens evolve after firms exhaust inventories or reorganize supply chains.

\par\textbf{Formulation and scope.} The source explicitly leaves aggregate price adjustment uncertain and suggests it may require more time. The whole-chain dynamic pass-through estimands, inventories and welfare accounting are proposed extensions; many short-run subcases are already estimated.

\par\textbf{Open-question provenance.} An explicit question or research agenda was verified at the source locator. This does not by itself certify that every later version remains unresolved \sref{30007624}{1}.

\par\textbf{Historical priority.} Earliest explicit uncertainty discussion located in this audit, not the historical first statement of tariff-incidence research. The March 2026 conference manuscript predates its April NBER issue.

\subsection{Short English statement}
In the explicitly specified setting: How does the burden of large tariffs shift over time among exporters, importers, retailers, workers and consumers? The mathematical statement above fixes the target, assumptions and scope of this catalogue formulation.

\subsection{Sources}
\begin{itemize}\raggedright
\item[\textnormal{[S1]}]\hypertarget{source:30007624:1}{} Pablo D. Fajgelbaum, Amit Khandelwal. 2026. Tariffs in 2025: Short-Run Impacts on the U.S. Economy. Abstract, PDF page 2; Section 5.2.2, final paragraph, PDF page 30; Section 6.4, PDF page 44.
\url{https://www.brookings.edu/wp-content/uploads/2026/03/1_Fajgelbaum-Khandelwal_unembargoed.pdf}.
DOI: 10.3386/w35064.
{\small\textit{Earliest verified question/agenda reference; historical priority qualified below. Finds terms-of-trade evidence inconclusive; longer adjustment and manufacturing outcomes remain uncertain.}}
\end{itemize}
\end{document}
